\chapter{Funkcje monotoniczne}
\label{ma1:ch:monotone-functions}

Porządek liczb rzeczywistych nakłada na funkcje monotoniczne silne ograniczenia. Nawet bez założenia ciągłości można opisać ich zachowanie graniczne, a wszystkie punkty nieciągłości mają szczególną postać. W tym rozdziale dziedziny i zbiory wartości rozpatrujemy ze zwykłą metryką odziedziczoną z $\R$.

\section{Monotoniczność}
\begin{definicja}
\label{ma1:def:monotone-function}
Niech $A\subseteq\R$ i $f:A\to\R$. Funkcja $f$ jest niemalejąca, jeśli dla wszystkich $x,y\in A$ z $x<y$ zachodzi $f(x)\le f(y)$; jest nierosnąca, jeśli zachodzi $f(x)\ge f(y)$. Jeżeli w tych warunkach wymagamy odpowiednio $f(x)<f(y)$ lub $f(x)>f(y)$, mówimy, że $f$ jest ściśle rosnąca lub ściśle malejąca. Funkcję nazywamy monotoniczną, jeśli jest niemalejąca lub nierosnąca, a ściśle monotoniczną, jeśli jest ściśle rosnąca lub ściśle malejąca.
\end{definicja}
Przejście od $f$ do $-f$ zamienia funkcję niemalejącą na nierosnącą i odwrotnie; zachowuje też ścisłą monotoniczność. Funkcja stała jest zarówno niemalejąca, jak i nierosnąca. Na zbiorze zawierającym co najmniej dwa punkty nie jest ściśle monotoniczna.

\begin{przyklad}
\label{ma1:ex:monotone-basic}
Funkcja $x\mapsto x^2$ jest ściśle rosnąca na $[0,+\infty)$ i ściśle malejąca na $(-\infty,0]$. Wynika to z tożsamości $y^2-x^2=(y-x)(y+x)$ dla $x<y$. Nie jest monotoniczna na $\R$. Funkcja $x\mapsto\lfloor x\rfloor$ jest niemalejąca na $\R$: jeśli $x<y$, to $\lfloor x\rfloor$ jest liczbą całkowitą nie większą od $y$, zatem $\lfloor x\rfloor\le\lfloor y\rfloor$. Nie jest ściśle rosnąca, ponieważ na każdym przedziale $[k,k+1)$, $k\in\Z$, jest stała.
\end{przyklad}

\section{Granice jednostronne}
\begin{twierdzenie}
\label{ma1:thm:monotone-one-sided}
Niech $I\subseteq\R$ będzie niepustym przedziałem, $f:I\to\R$ funkcją niemalejącą, a $p\in\R$. Jeżeli $p$ jest punktem skupienia zbioru $I_-=I\cap(-\infty,p)$, to
\[
 f(p^-):=\lim_{I\ni x\to p^-}f(x)=\sup\{f(x):x\in I_-\}
 \quad\text{w }\R\cup\{+\infty\}.
\]
Jeżeli $p$ jest punktem skupienia zbioru $I_+=I\cap(p,+\infty)$, to
\[
 f(p^+):=\lim_{I\ni x\to p^+}f(x)=\inf\{f(x):x\in I_+\}
 \quad\text{w }\R\cup\{-\infty\}.
\]
Dla funkcji nierosnącej w pierwszym wzorze występuje kres dolny, a w drugim kres górny, z odpowiednio zamienionymi dopuszczalnymi granicami nieskończonymi.
\end{twierdzenie}
\begin{proof}
Załóżmy najpierw, że $p\in I_-'$. Zbiór wartości na $I_-$ jest niepusty. Niech $s$ będzie jego kresem górnym w rozszerzonej prostej rzeczywistej. Jeżeli $s\in\R$, to dla $\varepsilon>0$ istnieje $u\in I_-$ z $s-\varepsilon<f(u)\le s$. Przyjmijmy $\delta=p-u>0$. Dla $x\in I$ z $p-\delta<x<p$ mamy $u<x$, stąd
$s-\varepsilon<f(u)\le f(x)\le s$.
To daje granicę $s$. Jeżeli $s=+\infty$, to dla dowolnego $M>0$ wybieramy $u\in I_-$ z $f(u)>M$. Ten sam wybór $\delta=p-u$ daje $f(x)\ge f(u)>M$ w odpowiednim lewostronnym otoczeniu punktu $p$.

Dla granicy prawostronnej niech $l$ będzie kresem dolnym wartości na $I_+$. Jeśli $l\in\R$, wybieramy $v\in I_+$ z $l\le f(v)<l+\varepsilon$. Dla $p<x<p+\delta$, gdzie $\delta=v-p$, i $x\in I$ mamy $l\le f(x)\le f(v)<l+\varepsilon$. Jeśli $l=-\infty$, dla każdego $M>0$ wybieramy $v$ z $f(v)<-M$ i analogicznie otrzymujemy $f(x)<-M$. Dla funkcji nierosnącej stosujemy otrzymane wzory do $-f$, korzystając z tego, że negacja zamienia kres górny na dolny i odwrotnie.
\end{proof}

\begin{wniosek}
\label{ma1:cor:monotone-interior-limits}
Niech $f:I\to\R$ będzie funkcją niemalejącą na przedziale $I$. Dla każdego punktu wewnętrznego $p\in I$ istnieją skończone granice jednostronne oraz
$f(p^-)\le f(p)\le f(p^+)$. Funkcja jest ciągła w $p$ wtedy i tylko wtedy, gdy $f(p^-)=f(p^+)$. Dla funkcji nierosnącej nierówności mają przeciwny zwrot, a kryterium ciągłości jest takie samo.
\end{wniosek}
\begin{proof}
Obie części przedziału po lewej i prawej stronie $p$ są niepuste i mają punkt skupienia $p$. Monotoniczność daje $f(x)\le f(p)\le f(y)$ dla $x<p<y$ w $I$. Zatem kres górny po lewej jest skończony i nie większy od $f(p)$, a kres dolny po prawej jest skończony i nie mniejszy od $f(p)$. Jeżeli granice są równe, nierówności wymuszają ich równość z $f(p)$, więc z twierdzenia \ref{ma1:thm:one-sided-limit-equivalence} i definicji ciągłości otrzymujemy tezę. Odwrotnie, ciągłość implikuje obie granice równe $f(p)$. Przypadek funkcji nierosnącej redukujemy do $-f$.
\end{proof}
W punkcie końcowym przedziału należącym do dziedziny ciągłość sprawdzamy względem dziedziny: jeśli $a=\min I$ i $I$ zawiera punkt większy od $a$, to dla funkcji niemalejącej granica $f(a^+)$ jest skończona (jest ograniczona od dołu przez $f(a)$ i od góry przez dowolną wartość po prawej). Ciągłość w $a$ jest równoważna $f(a^+)=f(a)$. Analogicznie w prawym końcu, a dla funkcji nierosnących po zmianie znaku. Na przedziale jednoelementowym każda funkcja jest ciągła.

\begin{twierdzenie}
\label{ma1:thm:monotone-at-infinity}
Niech $I\subseteq\R$ będzie niepustym przedziałem, a $f:I\to\R$ funkcją niemalejącą. Jeśli $I$ jest nieograniczony z góry, to $\lim_{I\ni x\to+\infty}f(x)=\sup f(I)$ w $\R\cup\{+\infty\}$. Jeśli $I$ jest nieograniczony z dołu, to $\lim_{I\ni x\to-\infty}f(x)=\inf f(I)$ w $\R\cup\{-\infty\}$. Dla funkcji nierosnącej kresy i dopuszczalne nieskończoności zamieniają się rolami.
\end{twierdzenie}
\begin{proof}
Jeśli $s=\sup f(I)$ jest skończony, dla $\varepsilon>0$ wybieramy $u\in I$ z $s-\varepsilon<f(u)\le s$. Dla $x\in I$ z $x>\max(1,u)$ zachodzi $s-\varepsilon<f(x)\le s$. Jeśli $s=+\infty$, dla $M>0$ wybieramy $u$ z $f(u)>M$ i otrzymujemy $f(x)>M$ dla $x>\max(1,u)$. Dla granicy przy $-\infty$ stosujemy właśnie udowodnioną część do funkcji $g(t)=-f(-t)$ na przedziale $-I$, nieograniczonym z góry. Funkcja $g$ jest niemalejąca, a $\sup g(-I)=-\inf f(I)$. Zmiana znaku argumentu i wartości w definicji granicy daje żądaną tezę. Przypadek nierosnący wynika z rozważenia $-f$.
\end{proof}

\section{Punkty nieciągłości i skoki}
\begin{definicja}
\label{ma1:def:discontinuity-types}
Niech $I$ będzie przedziałem, $p$ jego punktem wewnętrznym i $f:I\to\R$. Jeżeli obie granice jednostronne w $p$ istnieją i są skończone, lecz $f$ nie jest ciągła w $p$, mówimy o nieciągłości pierwszego rodzaju. Jeżeli granice te są równe, nieciągłość nazywamy usuwalną; jeżeli są różne, nieciągłością skokową. Pozostałe punkty nieciągłości wewnątrz $I$ nazywamy punktami nieciągłości drugiego rodzaju.
\end{definicja}
W punkcie nieciągłości usuwalnej zmiana wartości $f(p)$ na wspólną granicę jednostronną daje ciągłość w $p$. W punkcie skokowym żadna zmiana samej wartości $f(p)$ nie wystarcza, gdyż nie wpływa na żadną granicę jednostronną. Dla funkcji monotonicznej wszystkie nieciągłości w punktach wewnętrznych są skokowe: wynika to bezpośrednio z wniosku \ref{ma1:cor:monotone-interior-limits}. Dla funkcji niemalejącej rozmiar skoku w punkcie wewnętrznym oznaczamy przez $J_f(p)=f(p^+)-f(p^-)\ge0$.

\begin{przyklad}
\label{ma1:ex:monotone-jump}
Niech $0\le c\le1$ i $f:\R\to\R$ będzie określona przez $f(x)=0$ dla $x<0$, $f(0)=c$, $f(x)=1$ dla $x>0$. Funkcja jest niemalejąca, w zerze ma granice $0$ i $1$ oraz skok $J_f(0)=1$ niezależnie od wyboru $c$. Dla $c\notin[0,1]$ monotoniczność zostaje naruszona: porównanie wartości w zerze z wartościami po lewej lub prawej stronie daje sprzeczność.
\end{przyklad}
\begin{center}
\begin{tikzpicture}[x=2cm,y=2cm,>=Stealth]
 \draw[->] (-1.6,0)--(1.6,0) node[right] {$x$};
 \draw[->] (0,-.22)--(0,1.35) node[above] {$y$};
 \draw[very thick,blue!65!black] (-1.45,0)--(0,0);
 \draw[very thick,blue!65!black] (0,1)--(1.45,1);
 \draw[thick,blue!65!black,fill=white] (0,0) circle (2.2pt);
 \draw[thick,blue!65!black,fill=white] (0,1) circle (2.2pt);
 \fill[blue!65!black] (0,.4) circle (2pt) node[left] {$c$};
 \node[below right] at (0,0) {$0$};
 \node[above left] at (0,1) {$1$};
 \draw[<->] (.38,0)--(.38,1) node[midway,right] {$J_f(0)$};
\end{tikzpicture}

\small Rysunek: skok funkcji niemalejącej przy $0<c<1$; wartość w punkcie leży między granicami jednostronnymi.
\end{center}

\begin{lemat}
\label{ma1:lem:disjoint-jump-intervals}
Niech $f:I\to\R$ będzie niemalejąca, a $p<q$ będą punktami wewnętrznymi przedziału $I$. Wówczas $f(p^+)\le f(q^-)$.
\end{lemat}
\begin{proof}
Punkt $r=(p+q)/2$ należy do $I$. Z charakterystyki kresów w twierdzeniu \ref{ma1:thm:monotone-one-sided} mamy $f(p^+)\le f(r)\le f(q^-)$. W szczególności otwarte przedziały $(f(p^-),f(p^+))$ i $(f(q^-),f(q^+))$ są rozłączne.
\end{proof}

\begin{twierdzenie}
\label{ma1:thm:monotone-countable-discontinuities}
Niech $I$ będzie przedziałem i $f:I\to\R$ funkcją monotoniczną. Zbiór punktów nieciągłości $f$ względem $I$ jest co najwyżej przeliczalny.
\end{twierdzenie}
\begin{proof}
Wystarczy przypadek funkcji niemalejącej, ponieważ $f$ i $-f$ mają te same punkty ciągłości. Niech $D$ będzie zbiorem punktów nieciągłości wewnątrz $I$. Dla każdego $p\in D$ mamy $f(p^-)<f(p^+)$. Ustalmy ciąg $(r_n)_{n\ge1}$ zawierający wszystkie liczby wymierne. Z ich gęstości istnieje indeks $n$ z $f(p^-)<r_n<f(p^+)$. Przypiszmy punktowi $p$ najmniejszy taki indeks $n(p)$. Przedziały odpowiadające różnym punktom są rozłączne z lematu \ref{ma1:lem:disjoint-jump-intervals}, więc $n(p)\ne n(q)$ dla $p\ne q$. Otrzymujemy iniekcję $D\to\N$. Punkty $I$ niebędące jego punktami wewnętrznymi to co najwyżej dwa końce przedziału. Dodanie co najwyżej dwóch punktów do zbioru co najwyżej przeliczalnego zachowuje tę własność.
\end{proof}

\begin{twierdzenie}
\label{ma1:thm:finite-large-jumps}
Niech $a,b\in\R$, $a<b$, $f:[a,b]\to\R$ będzie niemalejąca i $\eta>0$. Zbiór punktów $p\in(a,b)$, dla których $J_f(p)\ge\eta$, jest skończony, a jego liczność nie przekracza $(f(b)-f(a))/\eta$.
\end{twierdzenie}
\begin{proof}
Dla dowolnych różnych $p_1<\cdots<p_m$ należących do tego zbioru oznaczmy $l_i=f(p_i^-)$ i $u_i=f(p_i^+)$. Z monotoniczności i lematu \ref{ma1:lem:disjoint-jump-intervals} mamy
$f(a)\le l_1\le u_1\le l_2\le u_2\le\cdots\le l_m\le u_m\le f(b)$.
Stąd
\[
 m\eta\le\sum_{i=1}^m(u_i-l_i)
 =u_m-l_1-\sum_{i=1}^{m-1}(l_{i+1}-u_i)
 \le f(b)-f(a).
\]
Gdyby zbiór był nieskończony, można byłoby wybrać $m>(f(b)-f(a))/\eta$ różnych punktów, co przeczyłoby nierówności.
\end{proof}

\begin{przyklad}
\label{ma1:ex:dense-monotone-jumps}
Zbiór nieciągłości funkcji monotonicznej może być gęsty. Niech $(r_n)_{n\ge1}$ będzie wyliczeniem bez powtórzeń zbioru $\Q\cap(0,1)$ i
\[
 f(x)=\sum_{n=1}^\infty 2^{-n}\chi_{[r_n,1]}(x),
 \qquad x\in[0,1],
\]
gdzie $\chi_E(x)=1$ dla $x\in E$ i $0$ poza $E$. Szereg istnieje w każdym punkcie przez porównanie z szeregiem geometrycznym. Każdy składnik jest niemalejący, a przejście do granicy sum częściowych zachowuje nierówność, więc $f$ jest niemalejąca.

Dla $p=r_k$ i $x<p$ mamy $f(p)-f(x)\ge2^{-k}$, ponieważ składnik o indeksie $k$ zwiększa się o tę wartość, a pozostałe nie maleją. Zatem $f(p^-)\le f(p)-2^{-k}$ i $f$ nie jest ciągła w $p$.

Jeśli $p\in[0,1]\setminus\{r_n:n\ge1\}$, dla $\varepsilon>0$ wybieramy $N\ge1$ z $2^{-N}<\varepsilon$. Niech $\delta=\min_{1\le n\le N}|p-r_n|>0$. Dla $x\in[0,1]$ z $|x-p|<\delta$ pierwsze $N$ składników przyjmuje w $x$ i $p$ te same wartości. Pozostała różnica jest ograniczona przez $\sum_{n>N}2^{-n}=2^{-N}<\varepsilon$. Dowodzi to ciągłości również w końcach przedziału. Zbiór nieciągłości jest dokładnie równy $\Q\cap(0,1)$.
\end{przyklad}

\section{Monotoniczność a własność Darboux}
\begin{definicja}
\label{ma1:def:darboux-property}
Niech $I$ będzie przedziałem. Mówimy, że funkcja $f:I\to\R$ ma własność Darboux, jeżeli dla wszystkich $x,y\in I$ z $x<y$ oraz każdej liczby $t$ między $f(x)$ i $f(y)$, z końcami włącznie, istnieje $z\in[x,y]$ z $f(z)=t$.
\end{definicja}
\begin{twierdzenie}
\label{ma1:thm:monotone-darboux-continuity}
Niech $I$ będzie przedziałem i $f:I\to\R$ funkcją monotoniczną. Wówczas $f$ jest ciągła na $I$ wtedy i tylko wtedy, gdy ma własność Darboux.
\end{twierdzenie}
\begin{proof}
Ciągłość implikuje własność Darboux przez zastosowanie twierdzenia \ref{ma1:thm:intermediate-value} na każdym $[x,y]\subseteq I$. Odwrotnie, wystarczy przypadek niemalejący, gdyż negacja zachowuje własność Darboux.

Załóżmy, że $p$ jest punktem nieciągłości. Jeżeli istnieje lewa strona z punktem skupienia $p$ i $f(p^-)<f(p)$, wybieramy $t$ z $f(p^-)<t<f(p)$ oraz $x\in I$ z $x<p$. Na $[x,p)$ wszystkie wartości są nie większe od $f(p^-)$, a w $p$ wartość jest większa od $t$. Własność Darboux dla $x,p$ wymaga jednak przyjęcia wartości $t$ na $[x,p]$, sprzeczność. Jeżeli istnieje prawa strona z punktem skupienia $p$ i $f(p)<f(p^+)$, wybieramy $t$ z $f(p)<t<f(p^+)$ oraz $y>p$ w $I$. Na $(p,y]$ wszystkie wartości są nie mniejsze od $f(p^+)$, więc wartość $t$ nie jest osiągana na $[p,y]$, również sprzecznie z własnością Darboux. W punkcie wewnętrznym nieciągłość wymusza co najmniej jedną z tych ścisłych nierówności; w końcu należącym do $I$ wymusza nierówność po istniejącej stronie. Przedział jednoelementowy nie ma punktów nieciągłości. Otrzymujemy ciągłość na całym $I$.
\end{proof}

\section{Funkcje odwrotne}
\begin{twierdzenie}
\label{ma1:thm:strict-monotone-inverse}
Niech $I$ będzie niepustym przedziałem i $f:I\to\R$ funkcją ściśle monotoniczną. Wówczas $f$ jest bijekcją na $f(I)$, jej odwrotność $g:f(I)\to I$ jest ściśle monotoniczna o tym samym kierunku monotoniczności i jest ciągła względem $f(I)$, nawet jeśli $f$ nie jest ciągła.
\end{twierdzenie}
\begin{proof}
Ścisła monotoniczność daje iniektywność, a surjektywność na obraz wynika z jego definicji. Dla $u<v$ w $f(I)$, jeśli $f$ jest ściśle rosnąca, to $g(u)<g(v)$: przeciwna nierówność implikowałaby $u\ge v$. Dla funkcji ściśle malejącej otrzymujemy $g(u)>g(v)$.

Pokażmy ciągłość $g$ dla ściśle rosnącego $f$. Niech $u_0=f(p)$, $p\in I$ i $\varepsilon>0$. Jeżeli istnieją punkty $x\in I$ z $x<p$, wybieramy $a\in I$ z $p-\varepsilon<a<p$: dla dowolnego $x<p$ w $I$ można wziąć $a=p-\min(\varepsilon,p-x)/2$, który należy do $[x,p]\subseteq I$. Wtedy $f(a)<u_0$. Jeśli istnieją punkty $y>p$ w $I$, analogicznie wybieramy $b\in I$ z $p<b<p+\varepsilon$, co daje $f(b)>u_0$.

Niech $\delta>0$ będzie minimum liczby $1$ oraz wszystkich istniejących dodatnich różnic $u_0-f(a)$ i $f(b)-u_0$. Jeśli $u\in f(I)$ i $|u-u_0|<\delta$, to gdy $a$ istnieje, mamy $u>f(a)$, więc $g(u)>a>p-\varepsilon$. Gdy $a$ nie istnieje, wszystkie punkty $I$ są nie mniejsze od $p$, co również daje $g(u)>p-\varepsilon$. Analogicznie, gdy $b$ istnieje, mamy $u<f(b)$ i $g(u)<b<p+\varepsilon$, a w przeciwnym przypadku $g(u)\le p<p+\varepsilon$. Zatem $|g(u)-g(u_0)|<\varepsilon$. Dla funkcji malejącej stosujemy wynik do $-f$; jej odwrotność jest równa $u\mapsto g(-u)$ na $-f(I)$, a złożenie z ciągłą negacją daje ciągłość $g$.
\end{proof}

\begin{przyklad}
\label{ma1:ex:inverse-of-jump}
Funkcja $f:\R\to\R$ dana przez $f(x)=x$ dla $x<0$ oraz $f(x)=x+1$ dla $x\ge0$ jest ściśle rosnąca. Nie jest ciągła w zerze, gdzie $f(0^-)=0$ i $f(0^+)=1$. Jej obraz to $B=(-\infty,0)\cup[1,+\infty)$. Odwrotność $g:B\to\R$ ma postać $g(u)=u$ dla $u<0$ i $g(u)=u-1$ dla $u\ge1$. Jest ciągła względem $B$ na mocy twierdzenia \ref{ma1:thm:strict-monotone-inverse}. W szczególności w $u=1$ promień $\delta=\min(1,\varepsilon)$ wyklucza wszystkie argumenty ujemne, więc $|g(u)-g(1)|=|u-1|<\varepsilon$ dla $u\in B$ z $|u-1|<\delta$. Nie należy wymagać granicy lewostronnej względem $\R$ w punkcie, w którym dziedzina odwrotności nie ma lewostronnego punktu skupienia.
\end{przyklad}

\begin{twierdzenie}
\label{ma1:thm:continuous-injective-monotone}
Niech $I$ będzie niepustym przedziałem i $f:I\to\R$ ciągłą iniekcją. Wówczas $f$ jest ściśle monotoniczna.
\end{twierdzenie}
\begin{proof}
Dla $x<y<z$ w $I$ wartość $f(y)$ musi leżeć ściśle między $f(x)$ i $f(z)$. Wszystkie trzy wartości są różne z iniektywności. Gdyby $f(y)>\max(f(x),f(z))$, wybieramy $t$ między $\max(f(x),f(z))$ i $f(y)$. Z twierdzenia o wartości pośredniej istnieją dwa różne punkty, jeden w $(x,y)$ i drugi w $(y,z)$, w których $f$ przyjmuje $t$, sprzecznie z iniektywnością. Przypadek $f(y)<\min(f(x),f(z))$ redukuje się do tego samego argumentu dla $-f$.

Jeżeli $I$ ma mniej niż dwa punkty, teza jest spełniona w sposób pusty. W przeciwnym razie wybierzmy $a<b$ w $I$. Po ewentualnej zmianie znaku możemy założyć $f(a)<f(b)$. Dla $x<y$ w $I$ uporządkujmy różne punkty zbioru $\{a,b,x,y\}$ jako $t_1<\cdots<t_k$. Dla każdych trzech kolejnych punktów własność właśnie udowodniona oznacza, że różnice $f(t_{i+1})-f(t_i)$ i $f(t_{i+2})-f(t_{i+1})$ są tego samego znaku. Zatem wszystkie kolejne różnice mają ten sam znak (dla $k=2$ jest tylko jedna różnica). Ponieważ $a,b$ należą do tego uporządkowanego zbioru i $f(a)<f(b)$, znak jest dodatni. W szczególności $f(x)<f(y)$. Dowodzi to ścisłego wzrostu; po cofnięciu zmiany znaku otrzymujemy ścisły spadek w drugim przypadku.
\end{proof}

\begin{wniosek}
\label{ma1:cor:monotone-homeomorphism}
Każda ciągła iniekcja $f:I\to\R$ określona na niepustym przedziale $I$ jest homeomorfizmem $I\to f(I)$, a $f(I)$ jest przedziałem.
\end{wniosek}
\begin{proof}
Obraz jest przedziałem na mocy wniosku \ref{ma1:cor:continuous-interval-image}. Ścisła monotoniczność wynika z twierdzenia \ref{ma1:thm:continuous-injective-monotone}, a ciągłość odwrotności z twierdzenia \ref{ma1:thm:strict-monotone-inverse}. Stosujemy definicję homeomorfizmu.
\end{proof}

Przy konstrukcji funkcji elementarnych wykorzystamy w szczególności istnienie granic funkcji monotonicznych oraz ciągłość odwrotności funkcji ściśle monotonicznych na przedziałach.
